% a05_twoloop_v: proof of lem:a05:D45 (v), part 1: scaling, forest form, the twelve graphs, lemmas (input from a05_cnum_NEG.tex)
\paragraph{Proof of \cref{lem:a05:D45}(v).} The proof rescales to unit scalar mass, writes King's coefficient as a forest-subtracted sum over the twelve graphs of order $e^4$, and bounds that sum uniformly in the rescaled photon mass. The power-counting input of \cite{Rei88,Rei88b} enters at one mass configuration only (scalar and photon mass $1$); the light-photon limit is handled by elementary bounds.

\begin{lemma}[exact scaling]\label{lem:a05:SC}
Write $d_{\alpha\beta}(s;\lmu;\eps)$ for the coefficient at renormalised scalar mass squared $s$, photon mass $\lmu$ and spacing $\eps$, and $n\defeq\alpha/2+\beta$. For every $t>0$, $d_{\alpha\beta}(s;\lmu;\eps)=t^{2-n}\,d_{\alpha\beta}(s/t^2;\lmu/t;t\eps)$. Consequently
\[
y_{40}(s;\lmu;\eps)=s^{-1}\,Y\big(\lmu/\sqrt s,\ \eps\sqrt s\big),\qquad Y(\nu,a)\defeq y_{40}(1;\nu;a).
\]
\end{lemma}
\begin{proof}
Consider a connected graph with $L$ loops, $I$ lines and vertices $v$ of lattice degree $d_v$ ($1$ for the current vertex, $0$ for the seagull and the quartic vertex, $2-n_v$ for the $n_v$-photon vertex of the hopping term), at zero external momentum. Substitute $q=tq'$ in every loop momentum: $q\in\BZ_\eps$ iff $q'\in\BZ_{t\eps}$, $d^3q=t^3d^3q'$, and $\frac4{\eps^2}\sin^2\frac{\eps tq'_\mu}2=t^2\cdot\frac4{(t\eps)^2}\sin^2\frac{(t\eps)q'_\mu}2$, so $\hat q^2$ at spacing $\eps$ equals $t^2$ times $\hat{q'}^2$ at spacing $t\eps$. Each propagator therefore gives $t^{-2}$: $(\hat q^2+s)^{-1}=t^{-2}(\hat{q'}^2+s/t^2)^{-1}$ and $(\hat k^2+\lmu^2)^{-1}=t^{-2}(\hat{k'}^2+\lmu^2/t^2)^{-1}$. A vertex of the hopping term is $\eps^{n_v-2}w(\eps q)$ with $w$ a trigonometric polynomial, and $\eps^{n_v-2}w(\eps q)=t^{2-n_v}(t\eps)^{n_v-2}w\big((t\eps)q'\big)$, i.e.\ $t^{d_v}$ times the same vertex at spacing $t\eps$; the quartic vertex is a constant. Hence the amplitude obeys $A(s,\lmu,\eps)=t^{3L-2I+\sum_vd_v}A(s/t^2,\lmu/t,t\eps)$, and for a two-point graph of order $n$ the exponent is the superficial degree $3-\frac E2-n=2-n$ ($E=2$). A counterterm vertex $d_{\alpha'\beta'}$ of lower order $n'$ inserted on a line contributes, by induction on the order, the factor $t^{2-n'}$ and one more propagator $t^{-2}$, which leaves the exponent $2-n$ unchanged. $d_{\alpha\beta}$ is the zero-momentum value of such a sum \cite[(3.32)--(3.33)]{K1}, so the identity holds. Differentiate in $s$ at fixed $t$ and put $t=\sqrt s$: $y_{\alpha\beta}(s;\lmu;\eps)=s^{-n/2}y_{\alpha\beta}(1;\lmu/\sqrt s;\eps\sqrt s)$, and $n=2$ at $(4,0)$.
\end{proof}

\paragraph{Conventions at unit scalar mass.} For the rest of the proof the spacing is $a\in(0,1]$, $\BZ_a\defeq[-\pi/a,\pi/a)^3$, $\langle f\rangle_a\defeq\int_{\BZ_a}f\,\frac{d^3k}{(2\pi)^3}$, $\langle\langle\cdot\rangle\rangle_a$ for two loop momenta; $\hat e(k)\defeq\sum_\mu\frac4{a^2}\sin^2\frac{ak_\mu}2$, so that $\frac4{\pi^2}|k|^2\le\hat e(k)\le|k|^2$ on $\BZ_a$ (\cref{lem:a05:el}(c)) and $\hat e\le12/a^2$; $\chi(k)\defeq\hat e(k)+1$, $C(q)\defeq\chi(q)^{-1}$ (scalar mass $1$), $D_\nu(k)\defeq(\hat e(k)+\nu^2)^{-1}$ (photon mass $\nu>0$; the propagator is diagonal in the photon indices), $D_0\defeq1/\hat e$. In the per-component normalisation of (i)--(v) the current vertex carries $\tilde V_\mu(p,p')\defeq-\frac2a\sin\frac{a(p+p')_\mu}2$ and the phase $e^{-\mathrm iak_\mu/2}$, $k=p'-p$ the incoming photon momentum; the seagull is $e^{-\mathrm iak_\mu/2}\cos\frac{a(p+p')_\mu}2\,\delta_{\mu\nu}$; the three- and four-photon vertices $V_3$, $V_4$ of the hopping term are $a$ and $a^2$ times trigonometric polynomials of $a(p+p')$ and $ak$ (\cref{lem:a05:deriv}). Numerical factors of the Feynman rules play no role and are suppressed. Soft and hard momenta: $Q\defeq[-\frac12,\frac12]^3\subset\BZ_a$, $Q_s\defeq\{|k|\le\frac12\}$, hard $\defeq\BZ_a\setminus Q_s$. For hard $k$, $\hat e(k)\ge\pi^{-2}$, so $D_\nu\le(1+\pi^2)D_1$; for every $k$, $D_\nu\le D_0\le\frac{\pi^2}4|k|^{-2}$ on $\BZ_a$. In particular, for every $\nu>0$ and $a\le1$,
\begin{equation}\label{eq:a05:soft}
\sigma_1\defeq\langle\mathbf 1_{Q_s}D_\nu\rangle_a\le\tfrac1{16},\qquad\langle\mathbf 1_{Q_s}D_\nu^2|k|^2\rangle_a\le\tfrac{\pi^2}{64},\qquad\langle C^2\rangle_a\le\tfrac{\pi^2}{64},\qquad\langle\chi^{-3/2}\rangle_a\le C\big(1+\ln\tfrac1a\big),
\end{equation}
by $\frac{\pi^2}4\int_{|k|\le1/2}\frac{d^3k}{(2\pi)^3|k|^2}=\frac1{16}$, $(\frac{\pi^2}4)^2\int_{|k|\le1/2}\frac{d^3k}{(2\pi)^3|k|^2}=\frac{\pi^2}{64}$, $(\frac{\pi^2}4)^2\int_{\R^3}\frac{d^3k}{(2\pi)^3(|k|^2+\pi^2/4)^2}=\frac{\pi^2}{64}$ and $(\frac{\pi^2}4)^{3/2}\int_{|k|\le\sqrt3\pi/a}\frac{d^3k}{(2\pi)^3(|k|^2+\pi^2/4)^{3/2}}$. Throughout, $C$ denotes a constant depending only on $a_0$ (fixed in \cref{lem:a05:H}), never on $\nu$, $a$ or $s$; it may change from line to line.

\paragraph{Subtraction operators, degrees, chords.} For a one-particle-irreducible subgraph $\gamma$ of $\Gamma$ with external momenta $q^\gamma$ and internal momenta $k^\gamma$, the subtraction operator of order $\delta\in\{0,1\}$ acts on the integrand $F_\gamma(k^\gamma;q^\gamma)$ at fixed $k^\gamma$ \cite[Def.~3.1]{Rei88b}:
\[
t^0_\gamma F_\gamma=F_\gamma(k^\gamma;0),\qquad t^1_\gamma F_\gamma=F_\gamma(k^\gamma;0)+\sum_\rho P_\rho(q^\gamma)\,\partial_{q_\rho}F_\gamma(k^\gamma;0),\qquad P_\rho(q)=a^{-1}p_\rho(aq),
\]
with $p_\rho$ smooth and $2\pi$-periodic, $p_\rho(x)=x_\rho+O(|x|^2)$ (\cite[Def.~3.1, (3-2), and App.~A]{Rei88b}; the same $P_\rho$ is used throughout). These operators preserve periodicity and act as the Taylor operators do on the degrees. The degree of $\gamma$ is $\omega(\gamma)=3L(\gamma)-2I(\gamma)+\sum_{v\in\gamma}d_v=3-\frac{E(\gamma)}2-n(\gamma)$ ($E$ external legs, $n$ the order); it is the UV-divergence degree of \cite[(2-28)--(2-30)]{Rei88b} in three dimensions, computed with the lattice degrees of \cite[Defs.~2.2--2.3]{Rei88}, the lattice degree of every vertex function being its continuum dimension. The lines of $\Gamma$ are ordered with the photon lines first, and the chord set of every connected subgraph is the one determined by this order through the rule of \cite[(2-24)]{Rei88b}: the chords of $\Gamma$ carry the loop momenta, all other line momenta follow by conservation, the chords of a subgraph $\gamma$ are chords of $\Gamma$, and the internal momenta of $\gamma$ are its chord momenta, which carry no external momentum of $\gamma$ \cite[(2-16)]{Rei88b}.

\begin{lemma}[the forest form of King's coefficient]\label{lem:a05:forest}
On $\eps\Z^3$, at any spacing and masses, with signed combinatorial weights $c_\Gamma$ whose values play no role below,
\[
d_{40}=\sum_\Gamma c_\Gamma\,\langle\langle R'F_\Gamma\rangle\rangle,\qquad y_{40}=\partial_sd_{40}=-\sum_\Gamma\sum_{l}c_\Gamma\,\langle\langle R'F_{\Gamma,\mathrm{dbl}(l)}\rangle\rangle ,
\]
where $\Gamma$ runs over the twelve one-particle-irreducible two-point graphs of order $e^4$ at zero external momentum (listed next), $F_\Gamma$ is the integrand of the Feynman rules without counterterms, $R'\defeq1-\sum_\gamma t^{\omega(\gamma)}_\gamma$ with $\gamma$ running over the proper one-particle-irreducible subgraphs with $\omega(\gamma)\ge0$, $l$ runs over the scalar lines of $\Gamma$, and $F_{\Gamma,\mathrm{dbl}(l)}$ is $F_\Gamma$ with the propagator of $l$ squared.
\end{lemma}
\begin{proof}
King's scheme \cite[(3.32)--(3.33)]{K1} has scalar mass counterterms at the orders $2\le\alpha+2\beta\le4$, each fixed by the vanishing at zero momentum of the one-particle-irreducible two-point function of its order with the lower-order counterterms inserted, and no photon or vertex counterterm. Compare it with the scheme in which every one-particle-irreducible type of non-negative degree receives the counterterm $-t^{\omega}$ of its zero-momentum Taylor polynomial, computed recursively in the order; its amplitudes are the zone integrals of the forest-subtracted integrands \cite[Sect.~4]{Rei88b}, and since the forests of proper divergent subgraphs have at most one element here (property (iv) below), the forest operator is $R'$. The proper divergent subgraphs of the graphs of order $e^4$ are of three types: (a) the one-loop scalar self-energy of order $e^2$ (exchange graph, seagull tadpole; $\omega=1$); (c) the one-loop vacuum polarisation of order $e^2$ (bubble, seagull loop; $\omega=1$); (d) the one-loop $\phi\bar\phi A$ vertex of order $e^3$ (triangle, two swordfish graphs, $A^3$-tadpole; $\omega=0$). The two schemes agree on $\eps\Z^3$:
(a) the order-$1$ part $t^1_\gamma-t^0_\gamma$ of the exchange integrand $f(q;p)=\widehat{(p+q)}^2D(p-q)C(q)$ is $\sum_\rho P_\rho(p)\partial_{p_\rho}f(q;0)$, which is odd in the loop momentum $q$ ($D$, $C$ even), while the seagull-tadpole integrand $D(k)\cos(ap_\mu)$ has zero gradient at $p=0$; so the order-$1$ term integrates to zero identically in the remaining momenta, and the order-$0$ term is King's $d_{20}$-insertion.
(c) The one-loop polarisation $\Pi^{(20)}_{\mu\nu}(k)$ (bubble plus seagull loop) satisfies $\Pi^{(20)}_{\mu\nu}(0)=\pm\big\langle\partial_\nu\big[\tfrac2a\sin(aq_\mu)\,C(q)\big]\big\rangle_a=0$ for all $\mu,\nu$, the integral of a derivative of a periodic function (on the diagonal the seagull loop is the derivative of the current loop; off the diagonal the bubble alone is a derivative; this is the lattice Ward identity used after \eqref{eq:a05:Pi}), and $\nabla_k\Pi^{(20)}(0)=0$: writing $\Pi^{(20)}_{\mu\nu}(k)=e^{-\mathrm ia(k_\mu-k_\nu)/2}T_{\mu\nu}(k)$, $T$ is even in $k$ (substitute $q\mapsto-q$), so $\nabla T(0)=0$ and $\nabla_k\Pi^{(20)}(0)=-\frac{\mathrm ia}2(e_\mu-e_\nu)\Pi^{(20)}_{\mu\nu}(0)=0$. Hence $t^1_k\Pi^{(20)}=0$: the photon counterterm of the forest scheme vanishes, as King's does.
(d) The vertex counterterm is the value at zero momenta of the sum of the one-loop vertex graphs, and each vanishes by $q\mapsto-q$: in the triangle $\tilde V_\mu(q,q)=-\frac2a\sin(aq_\mu)$ is odd, the other factors are even and the two phases of the internal photon cancel; the two swordfish graphs sum to $\langle2\mathrm i\sin(aq_\mu)D(q)C(q)\rangle_a=0$; the $A^3$-tadpole is proportional to $\sin\frac{a(p+p')_\mu}2=0$ at $p=p'=0$. So neither scheme has a vertex counterterm.
Therefore the zero-momentum value of the order-$e^4$ two-point function with King's insertions is $\sum_\Gamma c_\Gamma\langle\langle R'F_\Gamma\rangle\rangle$. Differentiating in $s$: the operators $t_\gamma$ act on momenta at fixed masses, so $\partial_s$ commutes with $R'$, and $\partial_sC=-C^2$ doubles one scalar line at a time.
\end{proof}

\paragraph{The twelve graphs.} They are the one-particle-irreducible two-point placements of the vertex sets $V_1^4$, $V_2V_1^2$, $V_2^2$, $V_3V_1$, $V_4$; every other placement contains a bridge. With the names of the table in \cref{lem:a05:40}:
\begin{center}\small
\begin{tabular}{@{}cllll@{}}
& vertices & graph & proper subgraphs with $\omega\ge0$ (type; photon lines inside) & class; removed terms of \eqref{eq:a05:OS}\\\hline
$\Gamma_0$ & $V_1^4$ & rainbow & inner exchange (a; 1) & A; $l=$ inner line\\
$\Gamma_1$ & $V_1^4$ & exchange, bubble inserted & bubble (c; 0) & B; $l$ in the bubble\\
$\Gamma_2$ & $V_1^4$ & crossed & two triangles (d; 1), overlapping & A; coefficient $0$\\
$\Gamma_3$ & $V_2V_1^2$ & exchange, seagull tadpole on the inner line & seagull tadpole (a; 1) & A; none\\
$\Gamma_4$ & $V_2V_1^2$ & seagull tadpole, bubble inserted & bubble (c; 0) & B; $l$ in the bubble\\
$\Gamma_5$ & $V_2V_1^2$ & seagull at an end (two mirror graphs) & swordfish vertex (d; 1) & A; coefficient $0$\\
$\Gamma_6$ & $V_2V_1^2$ & exchange, seagull loop inserted & seagull loop (c; 0) & B; vanishes, see (0)\\
$\Gamma_7$ & $V_2V_1^2$ & seagull in the middle & two swordfish vertices (d; 1), sharing $V_2$ & A; coefficient $0$\\
$\Gamma_8$ & $V_3V_1$ & $V_3$-tadpole vertex (two orientations) & $A^3$-tadpole (d; 1), no scalar line & A; none\\
$\Gamma_9$ & $V_2^2$ & seagull tadpole, seagull loop inserted & seagull loop (c; 0) & B; vanishes, see (0)\\
$\Gamma_{10}$ & $V_2^2$ & melon & none & A; none\\
$\Gamma_{11}$ & $V_4$ & $V_4$ double tadpole & none & no scalar line
\end{tabular}
\end{center}
By inspection of the twelve graphs: (i) $\omega(\Gamma)=0$, and for every scalar line $l$, $\omega_{\mathrm{dbl}}(\Gamma)=-2$, where $\omega_{\mathrm{dbl}}(\gamma)\defeq\omega(\gamma)-2\cdot\mathbf 1[l\in\gamma]$ is the degree of the doubled graph (the doubled line is a two-point insertion of degree $0$); (ii) every proper one-particle-irreducible subgraph with $\omega\ge0$ is one-loop, of type (a), (c) or (d), and contains at most one photon line, which, the photon lines being first in the order, is that subgraph's own chord; the subgraphs of type (c) contain no photon line; (iii) no subgraph containing $l$ has $\omega_{\mathrm{dbl}}\ge0$; (iv) two proper divergent subgraphs of the same $\Gamma$ always share a line or a vertex (the triangles of $\Gamma_2$ overlap, the swordfish vertices of $\Gamma_7$ share the seagull), so the forests of proper subgraphs in the sense of \cite[Sect.~2]{Rei88b} (members nested or with no common lines and vertices) have at most one element; (v) in class A both photon lines are chords of $\Gamma$, with independent loop momenta $k_1,k_2$, and the scalar lines carry integer combinations of them; in class B one photon line is a chord with momentum $k$, the other photon line carries $\pm k$, and the second chord is a scalar line of the closed scalar loop, with momentum $q$.

By (iv), $R'=1-\sum_\gamma t_\gamma$ splits according to whether $\gamma$ contains the doubled line: for every $\Gamma$ and scalar line $l$, pointwise,
\begin{equation}\label{eq:a05:OS}
R'F_{\Gamma,\mathrm{dbl}(l)}=R_{\mathrm{dbl}}F_{\Gamma,\mathrm{dbl}(l)}-\sum_{\gamma\ni l,\ \omega(\gamma)\ge0}\big(t^{\omega(\gamma)}_\gamma F_{\gamma,\mathrm{dbl}(l)}\big)\,F_{\bar\Gamma_\gamma},
\qquad R_{\mathrm{dbl}}\defeq1-\sum_{\gamma\not\ni l,\ \omega(\gamma)\ge0}t^{\omega(\gamma)}_\gamma ,
\end{equation}
where $F_{\bar\Gamma_\gamma}$ is the integrand of the lines and vertices of $\Gamma$ outside $\gamma$. $R_{\mathrm{dbl}}$ is the forest formula of the doubled graph with the labels $\omega_{\mathrm{dbl}}$ on every one-particle-irreducible subgraph: by (iii) the subgraphs containing $l$, and $\Gamma$ itself, have negative label and are not subtracted. The subtracted terms (``removed terms'') are polynomials of degree $\le\omega(\gamma)$ in $P(q^\gamma)$ whose coefficients are integrals over $k^\gamma$.

\begin{lemma}[photon factorisation]\label{lem:a05:PF}
Fix $\Gamma$, a scalar line $l$ and the photon mass $\nu$, and write $F^{(\nu)}$ for the integrand at that photon mass. In class A, $R_{\mathrm{dbl}}F^{(\nu)}_{\Gamma,\mathrm{dbl}(l)}(k_1,k_2)=D_\nu(k_1)D_\nu(k_2)\,R_{\Gamma,l}(k_1,k_2)$, and each removed term of \eqref{eq:a05:OS} is $D_\nu(k^\gamma)D_\nu(k^{\bar\Gamma})$ times a $\nu$-free function. In class B, $R_{\mathrm{dbl}}F^{(\nu)}_{\Gamma,\mathrm{dbl}(l)}(k,q)=D_\nu(k)^2R_{\Gamma,l}(k,q)$, and the same holds for $R'F^{(\nu)}_{\mathrm{dbl}}$ and for the removed term. The functions $R_{\Gamma,l}$ and the cofactors do not depend on $\nu$.
\end{lemma}
\begin{proof}
Every factor of $F_\Gamma$ other than the two photon propagators is $\nu$-free: the vertex functions are mass-free trigonometric polynomials and the scalar lines have mass $1$. A subtraction operator $t_\gamma$ Taylor-expands the integrand of $\gamma$ in $\gamma$'s external momenta at fixed internal momenta, and $\gamma$'s internal momentum is its chord momentum, which carries no external momentum of $\gamma$ \cite[(2-16), (2-24)]{Rei88b}. By (ii) every subtracted $\gamma$ contains either exactly one photon line, which is $\gamma$'s own chord, so that its propagator is a function of $\gamma$'s internal momentum alone and is untouched by $t_\gamma$, or no photon line (type (c)), its photon legs being external to $\gamma$, so that the attached photon propagators belong to $\bar\Gamma_\gamma$ and are cofactors. Hence no subtraction operator acts on a photon propagator, and the photon propagators factor out of every term of \eqref{eq:a05:OS} with the momenta of (v).
\end{proof}

\begin{lemma}[the fixed-mass bound]\label{lem:a05:H}
There are $a_0\in(0,1]$ and $K<\infty$ such that, at scalar and photon mass $1$, for every $\Gamma$ and every scalar line $l$ of $\Gamma$,
\[
\sum_{|\beta|\le6}\ \sup_{0<a\le a_0}\ \big\langle\big\langle\big|\partial^\beta R_{\mathrm{dbl}}F^{(1)}_{\Gamma,\mathrm{dbl}(l)}\big|\big\rangle\big\rangle_a\le K,
\]
$\partial^\beta$ running over the derivatives in the two loop momenta.
\end{lemma}
\begin{proof}
$R_{\mathrm{dbl}}F^{(1)}$ is a forest-subtracted lattice Feynman integrand in the sense of \cite{Rei88b}: the propagators are of the form $(\hat e(l)+\mu^2)^{-1}$ with $\mu=1$, the numerators are trigonometric polynomials of the class of \cite[Defs.~2.1--2.5]{Rei88} whose lattice degrees equal their continuum dimensions, the doubled line is a two-point insertion of degree $0$, and the labels $\delta(\gamma)=\omega_{\mathrm{dbl}}(\gamma)$ are the degrees of the doubled graph itself, so the conditions (4-3) and (4-4) of \cite{Rei88b} hold (remark after (4-4) there), with $t^\delta=0$ for $\delta<0$. The chord assignment obeys \cite[(2-24)]{Rei88b}, and the line momenta are natural in the sense of \cite[Def.~3.1]{Rei88}: every subtracted subgraph is one-loop, so its internal momentum is, up to sign, the momentum of one of its lines, and the set of line momenta of a connected graph without cut lines is natural (\cite{Rei88}, remark after Def.~3.1). By the renormalization theorem of \cite[Sects.~4--5]{Rei88b}, whose proof applies verbatim in three dimensions (the dimension enters only through the degree count; cf.\ the remark before Def.~2.5 of \cite{Rei88}), the integrand $R_{\mathrm{dbl}}F^{(1)}$ is $\frac{2\pi}a$-periodic, lies in the class of \cite{Rei88} and has negative degree on every Zimmermann subspace, the overall degree being $\omega_{\mathrm{dbl}}(\Gamma)=-2$; by \cite[Lemmas~2.1--2.2]{Rei88} the derivatives $\partial^\beta$ stay in the class and do not raise any of these degrees. The three-dimensional form of the renormalization theorem is read off its proof as follows. The dimension $d$ enters that proof at three places only: through the term $d\,m(\gamma)$, $m(\gamma)$ the number of loops, in the divergence degree \cite[(2-28), (2-31)]{Rei88b}; through the quantities $M_U(\gamma)=d\sum m$ of \cite[(5-10)]{Rei88b} and the bound $M_U(\Gamma)\ge dh$ of \cite[App.~C]{Rei88b}, which is $d$ times an inequality between loop counts and ranks of the coefficient matrices of the loop-momentum parametrisation of the Zimmermann subspace, the same for every $d$; and through the target inequality \cite[(5-5)]{Rei88b}, $dh+\operatorname{degr}_tR_\Gamma<0$. The degree bounds (5-11)--(5-13) of \cite[Lemma~5.2]{Rei88b} are proved by induction over the forest from the lattice degrees of the single propagators and vertex functions, which do not involve the dimension, from the subtraction degrees and from the loop counts, which add under the reduction of a subgraph, the same way as the degrees of (2-28) do; so with the degrees counted in $d$ dimensions and labels satisfying (4-3) and (4-4) in those degrees, every line of the induction holds for every $d$, and its conclusion in $d=3$, $3h+\operatorname{degr}_tR_\Gamma<0$ on every Zimmermann subspace, is the hypothesis (3.9) of the power counting theorem of \cite{Rei88}. Everything else in the proof (the subtraction operators and their degree calculus, the forest algebra, the periodicity of $R_\Gamma$) treats the momenta as lists of real variables and does not depend on the dimension. The power counting theorem \cite[Thm.~1]{Rei88} applies to each $\partial^\beta R_{\mathrm{dbl}}F^{(1)}$, and its proof bounds the integrand on every sector by non-negative majorants whose integrals over the zone are bounded uniformly in $a\le a_0$ when all degrees are negative (\cite[Sect.~4 and Thm.~2]{Rei88}), which is the uniform absolute bound claimed; $a_0$ and $K$ depend on the masses, here $(1,1)$, and on the finitely many integrands $(\Gamma,l,\beta)$.
\end{proof}
